\subsection{Mode count versus aperture and frequency}
\label{sec:ch481}
% TOC tag: [HOME]

Mode-count laws connect aperture area, wavelength, and the integer budgets used in array and MIMO design \cite{collin2001,jackson1999,hansen1988}.
Subsection~\ref{sec:ch481} restates those laws as upper bounds on \(d_{\mathrm{eff}}(\varepsilon)\) \cite{stratton1941,harrington2001,devaney2004}.

\paragraph{Assumptions.}
Presuppose Eqs.~\eqref{eq:ch4.truncation-volume-scaling}, \eqref{eq:ch4.truncation-area-scaling}, and Definition~\ref{def:ch4.effective-em-dimension}
\cite{stratton1941,hansen1988,harrington2001}.

\paragraph{Mode-count bounds.}
For ball \(B_{a}\),
\begin{equation}
\label{eq:ch4.mode-count-volume}
N_{\mathrm{mode}}^{\mathrm{(vol)}}(\varepsilon)
\ \le\
C_{V}(\varepsilon)\,(k_{0}a)^{3},
\end{equation}
and for aperture \(\Sigma_{a}\),
\begin{equation}
\label{eq:ch4.mode-count-area}
N_{\mathrm{mode}}^{\mathrm{(ap)}}(\varepsilon)
\ \le\
C_{A}(\varepsilon)\,k_{0}^{2}A_{a},
\end{equation}
with \(N_{\mathrm{mode}}\le d_{\mathrm{eff}}(\varepsilon)\) in the aligned export gauge \cite{collin2001,jackson1999,hansen1988,harrington2001}.
Equations~\eqref{eq:ch4.mode-count-volume}--\eqref{eq:ch4.mode-count-area} are realizability ceilings, not design targets
\cite{stratton1941,devaney2004}.

\begin{theorem}[Mode-count monotonicity]
\label{thm:ch4.mode-count-monotone}
At fixed \(\varepsilon\), \(N_{\mathrm{mode}}^{\mathrm{(vol)}}\) and \(N_{\mathrm{mode}}^{\mathrm{(ap)}}\) are non-decreasing in
\(k_{0}\) and in \(a\) or \(A_{a}\) until material dispersion alters \(\mathcal{J}_{\mathrm{real}}\) \cite{stratton1941,jackson1999,hansen1988}.
Phase-only tuning on fixed geometry does not raise \(N_{\mathrm{mode}}\) \cite{devaney2004,harrington2001}.
\end{theorem}

\paragraph{Worked illustration (5G panel).}
A \(k_{0}W\approx3\) panel has \(N_{\mathrm{mode}}^{\mathrm{(ap)}}\sim\mathcal{O}(10)\)--\(\mathcal{O}(30)\), not one per patch
\cite{collin2001,harrington2001,hansen1988}.

\paragraph{Problem (element count versus modes).}
\textbf{Problem.} \(64\) patches claim \(64\) independent beams. Audit.

\textbf{Step 1.} Compute \(k_{0}^{2}A_{a}\) \cite{jackson1999,stratton1941}.

\textbf{Step 2.} Bound \(N_{\mathrm{mode}}^{\mathrm{(ap)}}\) via Eq.~\eqref{eq:ch4.mode-count-area} \cite{hansen1988,harrington2001}.

\textbf{Step 3.} Report \(d_{\mathrm{eff}}(\varepsilon)\) from SVD \cite{devaney2004}.

\textbf{Physical meaning.} Elements sample the aperture; modes are geometry-bounded \cite{collin2001}.

\paragraph{Validity.}
Theorem~\ref{thm:ch4.mode-count-monotone} is asymptotic for \(k_{0}D_{s}\gtrsim1\) \cite{jackson1999,hansen1988}.


\subsection{High- and low-frequency asymptotics}
\label{sec:ch482}
% TOC tag: [HOME]

Radiative rank crosses distinct asymptotic regimes as \(k_{0}D_{s}\) varies \cite{stratton1941,jackson1999,hansen1988}. Subsection
~\ref{sec:ch482} records the low-frequency dipole limit and high-frequency volume-growth limit on \(d_{\mathrm{eff}}\)
\cite{harrington2001,devaney2004}.

\paragraph{Assumptions.}
Presuppose Eq.~\eqref{eq:ch4.coupling-crossover} and compact \(\Omega_{s}\) \cite{stratton1941,hansen1988}.

\paragraph{Rayleigh limit.}
As \(k_{0}D_{s}\to0\),
\begin{equation}
\label{eq:ch4.rayleigh-rank-limit}
d_{\mathrm{eff}}(\varepsilon)
\ \longrightarrow\
d_{0}\le3
\end{equation}
in the electric dipole class on centered \(\Omega_{s}\) \cite{jackson1999,stratton1941,harrington2001}. High-\(\ell\) synthesis is
forbidden by Eq.~\eqref{eq:ch4.high-ell-decay} in this regime \cite{hansen1988,devaney2004}.

\paragraph{Large-body limit.}
As \(k_{0}D_{s}\to\infty\) on fixed shape,
\begin{equation}
\label{eq:ch4.large-body-rank-growth}
d_{\mathrm{eff}}(\varepsilon)
\ \sim\
C_{\infty}(\varepsilon)\,(k_{0}D_{s})^{3},
\end{equation}
importing Eq.~\eqref{eq:ch4.mode-count-volume} \cite{stratton1941,jackson1999,hansen1988}. Truncation remains mandatory at any
finite tolerance \cite{harrington2001}.

\begin{corollary}[Asymptotic synthesis regimes]
\label{cor:ch4.asymptotic-regimes}
Inverse design on electrically small hardware is dipole-class limited; on electrically large hardware is volume-growth limited with
mandatory truncation at \(\varepsilon\) \cite{stratton1941,jackson1999,hansen1988,harrington2001,devaney2004}.
\end{corollary}

\paragraph{Worked illustration (UHF versus mmWave).}
A UHF loop at \(k_{0}D_{s}=0.3\) has \(d_{\mathrm{eff}}\le3\); a mmWave array at \(k_{0}D_{s}=4\) accesses dozens of export modes
\cite{jackson1999,harrington2001,collin2001}.

\paragraph{Problem (frequency sweep rank).}
\textbf{Problem.} Rank is measured at 3\,GHz and extrapolated to 30\,GHz on fixed hardware. Valid?

\textbf{Step 1.} Identify regime via \(k_{0}D_{s}\) at each frequency \cite{jackson1999,stratton1941}.

\textbf{Step 2.} Apply Eq.~\eqref{eq:ch4.rayleigh-rank-limit} or \eqref{eq:ch4.large-body-rank-growth} \cite{hansen1988,harrington2001}.

\textbf{Step 3.} Recompute \(d_{\mathrm{eff}}(\varepsilon)\) per frequency \cite{devaney2004}.

\textbf{Physical meaning.} Asymptotic regime sets the scaling law \cite{stratton1941}.

\paragraph{Validity.}
Eqs.~\eqref{eq:ch4.rayleigh-rank-limit}--\eqref{eq:ch4.large-body-rank-growth} assume linear isotropic exterior \cite{stratton1941}.


\subsection{Spectral concentration and localization}
\label{sec:ch483}
% TOC tag: [HOME]

Finite support concentrates spatial energy and compresses spectral support; concentration inequalities make truncation the optimal
description of exportable information \cite{stratton1941,jackson1999,hansen1988}. Subsection~\ref{sec:ch483} states concentration on
\(\boldsymbol{\Gamma}_\omega\) \cite{harrington2001,devaney2004}.

\paragraph{Assumptions.}
Presuppose Eq.~\eqref{eq:ch4.spectral-compression-def} and Theorem~\ref{thm:ch4.spectral-compression-optimal}
\cite{stratton1941,harrington2001,hansen1988}.

\paragraph{Concentration ratio.}
Define
\begin{equation}
\label{eq:ch4.spectral-concentration-ratio}
\Lambda_{\mathrm{conc}}(\varepsilon)
=
\frac{\sum_{n=1}^{d_{\mathrm{eff}}(\varepsilon)}\sigma_{n}^{2}}
{\sum_{n\ge1}\sigma_{n}^{2}}
\ \ge\
1-\varepsilon^{2},
\end{equation}
with \(\{\sigma_{n}\}\) from Eq.~\eqref{eq:ch4.gamma-SVD} \cite{harrington2001,hansen1988}. High \(\Lambda_{\mathrm{conc}}\) means
export energy concentrates in the leading \(d_{\mathrm{eff}}\) directions \cite{stratton1941,devaney2004}.

\begin{proposition}[Truncation as concentration optimum]
\label{prop:ch4.truncation-concentration}
Rank-\(d_{\mathrm{eff}}(\varepsilon)\) truncation maximizes \(\Lambda_{\mathrm{conc}}(\varepsilon)\) among all rank-\(N\) approximations
of \(\boldsymbol{\Gamma}_\omega\) \cite{harrington2001,devaney2004,stratton1941}. Alternative sparsity priors that retain higher
listing rank without raising \(\sigma_{n}\) above threshold do not increase export power \cite{hansen1988}.
\end{proposition}

\paragraph{Worked illustration (slow tail).}
If \(\Lambda_{\mathrm{conc}}(10^{-2})=0.85\), \(15\%\) of export gain lies outside the leading \(d_{\mathrm{eff}}\) subspace: truncation
at \(d_{\mathrm{eff}}\) sacrifices certified power, not merely labels \cite{harrington2001,hansen1988}.

\paragraph{Problem (concentration audit).}
\textbf{Problem.} Given \(\{\sigma_{n}\}\), find minimal \(N\) with \(\Lambda_{\mathrm{conc}}\ge0.99\).

\textbf{Step 1.} Accumulate \(\sigma_{n}^{2}\) until ratio Eq.~\eqref{eq:ch4.spectral-concentration-ratio} exceeds \(0.99\) \cite{harrington2001}.

\textbf{Step 2.} Set \(N=d_{\mathrm{eff}}(\varepsilon')\) at implied \(\varepsilon'\) \cite{hansen1988,devaney2004}.

\textbf{Step 3.} Compare to \(\ell_{\mathrm{supp}}\) for consistency \cite{stratton1941}.

\textbf{Physical meaning.} Concentration quantifies how sharply geometry compresses export \cite{hansen1988}.

\paragraph{Validity.}
Eq.~\eqref{eq:ch4.spectral-concentration-ratio} depends on flux normalization of \(\mathcal{C}\) \cite{hansen1988,stratton1941}.


\subsection{Landau-Pollak-Slepian limits and n-widths}
\label{sec:ch484}
% TOC tag: [HOME]

Prolate concentration and Kolmogorov \(n\)-widths bound how many degrees of freedom can be simultaneously localized and bandlimited
\cite{hansen1988,stratton1941}. Subsection~\ref{sec:ch484} imports those bounds as realizability ceilings on \(d_{\mathrm{eff}}\)
\cite{harrington2001,devaney2004}.

\paragraph{Assumptions.}
Consider spatial domain \(D\) supporting \(\Jimp\) and spectral domain \(\Omega_{k}\) supporting \(\widetilde{\mathbf{A}}_{\mathrm{rad}}\)
\cite{stratton1941,jackson1999}. Presuppose Eq.~\eqref{eq:ch4.duality-bandlimit-scaling} \cite{hansen1988}.

\paragraph{\(n\)-width realizability bound.}
Let \(N_{\mathrm{LPS}}(\varepsilon,D,\Omega_{k})\) denote the Kolmogorov \(n\)-width of the concentration operator linking \(D\) to
\(\Omega_{k}\) \cite{hansen1988,stratton1941}. Then
\begin{equation}
\label{eq:ch4.lps-nwidth-bound}
d_{\mathrm{eff}}(\varepsilon)
\ \le\
N_{\mathrm{LPS}}(\varepsilon,D,\Omega_{k}),
\end{equation}
because export synthesis cannot exceed simultaneous spatial--spectral concentration \cite{harrington2001,devaney2004}. On ball
\(B_{a}\), \(N_{\mathrm{LPS}}\sim(k_{0}a)^{3}\) consistent with Eq.~\eqref{eq:ch4.mode-count-volume} \cite{jackson1999,hansen1988}.

\begin{theorem}[LPS bound on synthesis dimension]
\label{thm:ch4.lps-synthesis-bound}
\(\dim\mathcal{S}_{\mathrm{real}}(\mathcal{C})\le N_{\mathrm{LPS}}(\varepsilon,D,\Omega_{k})\) for any admissible listing \(\mathcal{C}\)
\cite{stratton1941,hansen1988,harrington2001,devaney2004}.
\end{theorem}
\begin{proof}
Concentration operator dominates export map in operator norm on compact domains \cite{hansen1988}; Definition~\ref{def:ch4.effective-em-dimension}
\cite{harrington2001,stratton1941}.
\end{proof}

\paragraph{Worked illustration (prolate aperture).}
A prolate spheroidal aperture has anisotropic \(N_{\mathrm{LPS}}\): major-axis concentration raises \(d_{\mathrm{eff}}\) along one
direction before others \cite{hansen1988,jackson1999,harrington2001}.

\paragraph{Problem (LPS versus SVD rank).}
\textbf{Problem.} \(d_{\mathrm{eff}}(10^{-2})=40\) but \(N_{\mathrm{LPS}}=35\) for declared \((D,\Omega_{k})\). Which governs?

\textbf{Step 1.} Apply Theorem~\ref{thm:ch4.lps-synthesis-bound} as ceiling \cite{hansen1988,stratton1941}.

\textbf{Step 2.} If SVD exceeds LPS, audit domain declarations \cite{harrington2001,devaney2004}.

\textbf{Step 3.} Use \(\min(d_{\mathrm{eff}},N_{\mathrm{LPS}})\) for conservative design \cite{jackson1999}.

\textbf{Physical meaning.} \(n\)-widths are geometry--spectrum concentration limits \cite{hansen1988}.

\paragraph{Validity.}
Eq.~\eqref{eq:ch4.lps-nwidth-bound} assumes linear time-harmonic maps and declared concentration domains \cite{stratton1941,hansen1988}.

Section~\ref{sec:ch48} has linked mode count, asymptotics, concentration, and LPS bounds to \(d_{\mathrm{eff}}\)
\cite{stratton1941,hansen1988,harrington2001}. Section~\ref{sec:ch49} records physical consequences for engineering and measurement
\cite{devaney2004}.


\section{Physical Consequences of Realizability Limits}
\label{sec:ch49}

Truncation and rank are not abstract operator facts: they bound what finite apertures can radiate, what phase-only hardware can retune,
and what laboratories can observe \cite{stratton1941,harrington2001,devaney2004}. Section~\ref{sec:ch49} states those consequences in
design and measurement language without leaving the realizability framework established in Sections~\ref{sec:ch41}--\ref{sec:ch48}
\cite{hansen1988,jackson1999}.


\subsection{Geometry as origin of dimensionality limits}
\label{sec:ch491}
% TOC tag: [HOME]

Dimensionality limits in synthesis reports trace to \(\operatorname{diam}(\Omega_{s})\), \(A_{a}\), and \(k_{0}\) through
\(d_{\mathrm{eff}}(\varepsilon)\), not to solver tolerances alone \cite{stratton1941,jackson1999,hansen1988}. Subsection~\ref{sec:ch491}
states the origin principle for laboratory interpretation \cite{harrington2001,devaney2004}.

\paragraph{Assumptions.}
Presuppose Definition~\ref{def:ch4.effective-em-dimension} and Eqs.~\eqref{eq:ch4.mode-count-volume}--\eqref{eq:ch4.mode-count-area}
\cite{stratton1941,hansen1988,harrington2001}.

\begin{theorem}[Geometry as rank origin]
\label{thm:ch4.geometry-rank-origin}
At fixed \(\varepsilon\) and listing \(\mathcal{C}\), \(d_{\mathrm{eff}}(\varepsilon)\) is invariant under phase-only
reparameterizations of \(\Jimp\) on fixed \(\Omega_{s}\) and changes only when \(\operatorname{diam}(\Omega_{s})\), \(k_{0}\), or
\(\mathcal{J}_{\mathrm{real}}\) class changes \cite{stratton1941,jackson1999,hansen1988,harrington2001,devaney2004}.
\end{theorem}
\begin{proof}
\(\boldsymbol{\Gamma}_\omega\) depends on support and material class; phase rotations on fixed support are gauge transformations on
\(\mathcal{J}_{\mathrm{real}}\) that preserve singular values of the export map \cite{harrington2001,stratton1941}.
\end{proof}

\paragraph{Worked illustration (metasurface phase).}
A fixed-area metasurface that retunes phases without enlarging \(\mathcal{L}_{2}\) cannot raise \(d_{\mathrm{eff}}\)
\cite{stratton1941,harrington2001,hansen1988}.

\paragraph{Problem (solver rank versus hardware).}
\textbf{Problem.} FEM rank grows with mesh; diversity saturates. Explain.

\textbf{Step 1.} Report \(d_{\mathrm{eff}}(\varepsilon)\), not mesh DOF \cite{harrington2001,devaney2004}.

\textbf{Step 2.} Tie saturation to \(k_{0}D_{s}\) via Eq.~\eqref{eq:ch4.mode-count-volume} \cite{jackson1999,hansen1988}.

\textbf{Step 3.} Enlarge aperture or frequency to raise rank \cite{stratton1941}.

\textbf{Physical meaning.} Geometry is the origin of dimensionality \cite{stratton1941}.

\paragraph{Validity.}
Theorem~\ref{thm:ch4.geometry-rank-origin} assumes linear passive exterior \cite{stratton1941}.


\subsection{Limits of phase-only engineering}
\label{sec:ch492}
% TOC tag: [HOME]

Phase-only control explores a submanifold of \(\mathcal{J}_{\mathrm{real}}\) at fixed amplitude support \cite{stratton1941,harrington2001}.
Subsection~\ref{sec:ch492} states what phase-only designs can and cannot change in \(\mathcal{S}_{\mathrm{real}}(\mathcal{C})\)
\cite{devaney2004,hansen1988}.

\paragraph{Assumptions.}
Presuppose Proposition~\ref{prop:ch4.tangent-dimension} and Theorem~\ref{thm:ch4.geometry-rank-origin}
\cite{stratton1941,harrington2001,devaney2004}.

\begin{proposition}[Phase-only rank ceiling]
\label{prop:ch4.phase-only-ceiling}
Phase-only optimization on fixed \(\Omega_{s}\) cannot increase \(d_{\mathrm{eff}}(\varepsilon)\), \(\ell_{\mathrm{supp}}\), or
\(N_{\mathrm{mode}}\) \cite{stratton1941,jackson1999,hansen1988,harrington2001}. It rotates \(\mathbf{c}\) within the same export
subspace of dimension \(\le d_{\mathrm{eff}}(\varepsilon)\) \cite{devaney2004}.
\end{proposition}

\paragraph{Worked illustration (reflectarray).}
A reflectarray that steers a beam without enlarging aperture rotates \(\mathbf{c}\) within \(\ell=1\)--\(\ell_{\mathrm{supp}}\) content;
it does not manufacture high-\(\ell\) tails forbidden by Eq.~\eqref{eq:ch4.high-ell-decay} \cite{hansen1988,jackson1999,harrington2001}.

\paragraph{Problem (phase-only OAM claim).}
\textbf{Problem.} Phase plate claims new OAM order without size change. Audit.

\textbf{Step 1.} Test \(\ell_{\mathrm{supp}}\) before and after \cite{hansen1988,stratton1941}.

\textbf{Step 2.} Test tail exponent Eq.~\eqref{eq:ch4.tail-exponent} \cite{jackson1999,devaney2004}.

\textbf{Step 3.} Reject if only phase changed on fixed \(\mathcal{L}_{2}\) \cite{harrington2001}.

\textbf{Physical meaning.} Phase explores a fixed-rank manifold \cite{stratton1941}.

\paragraph{Validity.}
Proposition~\ref{prop:ch4.phase-only-ceiling} assumes passive linear exterior \cite{stratton1941}.


\subsection{Observable manifestations of truncation}
\label{sec:ch493}
% TOC tag: [USE]

Measurement reports from Section~\ref{sec:ch15} inherit truncation as smoothed patterns, elevated sidelobes, and saturated channel
counts \cite{stratton1941,harrington2001}. Subsection~\ref{sec:ch493} links those observables to \(\mathcal{B}_{\mathrm{fs}}\) and
\(d_{\mathrm{eff}}\) \cite{hansen1988,devaney2004}.

\paragraph{Assumptions.}
Presuppose Definition~\ref{def:ch4.finite-support-restriction} and Theorem~\ref{thm:ch4.finite-support-audit}
\cite{stratton1941,hansen1988,harrington2001}.

\paragraph{Observable signatures.}
Truncation manifests as
\begin{equation}
\label{eq:ch4.truncation-observables}
\Delta_{\mathrm{pat}}
\sim
\|\mathbf{R}_{\ell_{\max}}\|_{\mathrm{rad}},
\qquad
\Delta_{\mathrm{chan}}
\sim
d_{\mathrm{eff}}^{-1},
\end{equation}
with \(\mathbf{R}_{\ell_{\max}}\) from Eq.~\eqref{eq:ch3.truncation-remainder-energy} \cite{hansen1988,stratton1941}. Elevated
far-zone sidelobes and saturated MIMO capacity at fixed SNR are consistent signatures of \(\mathcal{A}_{\mathrm{fs}}>1\) or
\(d_{\mathrm{eff}}\) saturation \cite{harrington2001,devaney2004}.

\paragraph{Worked illustration (pattern shoulders).}
A target with plateau high-\(\ell\) tails produces shoulder artifacts when synthesized on \(B_{a}\): the artifact is a realizability
signature, not a calibration error \cite{hansen1988,jackson1999,harrington2001}.

\paragraph{Problem (diagnose truncation in data).}
\textbf{Problem.} Measured pattern shows unexpected shoulders. Truncation or miscalibration?

\textbf{Step 1.} Compute \(\mathcal{A}_{\mathrm{fs}}[\mathbf{c}]\) \cite{hansen1988,stratton1941}.

\textbf{Step 2.} Test tail mass Eq.~\eqref{eq:ch4.cumulative-tail-mass} \cite{harrington2001}.

\textbf{Step 3.} If audits fail, attribute to geometry rank before recalibrating \cite{devaney2004}.

\textbf{Physical meaning.} Truncation has characteristic far-zone signatures \cite{hansen1988}.

\paragraph{Validity.}
Eq.~\eqref{eq:ch4.truncation-observables} is narrowband \cite{stratton1941}.


\subsection{Implications for radiation and measurement}
\label{sec:ch494}
% TOC tag:

Laboratory certification requires \(\mathbf{c}\in\mathcal{S}_{\mathrm{real}}\cap\mathcal{B}_{\mathrm{fs}}\cap\mathcal{E}_{\mathrm{acc}}^{(N)}\)
from Theorem~\ref{thm:ch4.three-way-certificate} \cite{harrington2001,devaney2004,hansen1988}. Subsection~\ref{sec:ch494} closes the
measurement thread for Volume~I; full OAM tomography, MIMO capacity proofs, and inverse-source uniqueness on noisy data continue in
Volume~II \cite{stratton1941}.

\paragraph{Assumptions.}
Presuppose Eq.~\eqref{eq:ch4.certified-intersection-closure} and accessibility envelope
Definition~\ref{def:ch4.accessibility-envelope} \cite{harrington2001,devaney2004}.

\paragraph{Measurement closure in Volume I.}
Volume~I certifies that declared \((\Omega_{s},k_{0},\mathcal{C},N)\) yield a well-posed intersection
Eq.~\eqref{eq:ch4.certified-intersection-closure} \cite{stratton1941,hansen1988}. When \(\mathbf{c}\notin\mathcal{E}_{\mathrm{acc}}^{(N)}\),
the failure is instrument rank, not synthesis operator error \cite{harrington2001,devaney2004}. Broadband pulsed certification requires
\(\mathcal{B}_{\mathrm{fs}}(\omega)\) families not developed here \cite{jackson1999}.

\paragraph{Forward pointer.}
Volume~II develops measurement operators on multi-path channels, noise floors, and tomographic OAM readout beyond the narrowband
\(d_{\mathrm{eff}}\) certificate of Chapter~\ref{ch:04} \cite{harrington2001,devaney2004}.

\paragraph{Worked illustration (certified versus synthesizable).}
A synthesizable \(\mathbf{c}\) with \(\ell_{\mathrm{supp}}=8\) may fail certification at \(N=5\) modes: the laboratory reports a
rank-limited envelope, not a synthesis failure \cite{harrington2001,hansen1988}.

\paragraph{Problem (three-way failure diagnosis).}
\textbf{Problem.} Pattern error persists after calibration. Synthesis, finite support, or accessibility?

\textbf{Step 1.} Test \(\delta_{\mathrm{syn}}=0\) \cite{stratton1941,harrington2001}.

\textbf{Step 2.} Test \(\mathcal{A}_{\mathrm{fs}}\le1\) \cite{hansen1988}.

\textbf{Step 3.} Test \(\mathbf{c}\in\mathcal{E}_{\mathrm{acc}}^{(N)}\) \cite{devaney2004}.

\textbf{Physical meaning.} Measurement closes only on the three-way intersection \cite{stratton1941}.

\paragraph{Validity.}
Subsection~\ref{sec:ch494} does not develop Volume~II measurement protocols \cite{harrington2001}.

Section~\ref{sec:ch49} has identified geometry as rank origin, phase-only ceilings, observable truncation signatures, and Volume~I
measurement closure \cite{stratton1941,hansen1988,harrington2001}. Section~\ref{sec:ch410} supplies worked examples with figures
\cite{devaney2004}.


\section{Examples}
\label{sec:ch410}

The examples compress the realizability architecture into auditable calculations on finite apertures and feeds
\cite{stratton1941,harrington2001,hansen1988}. Each subsection states a problem, develops the governing inequalities, and interprets
the result in watts and export rank \cite{devaney2004}.


\subsection{Maximum angular momentum of a finite aperture}
\label{sec:ch4101}
% TOC tag: [FIG+PROB][HOME][USE SS3.7.2][USE SS3.4]

\textbf{Problem.} A circular aperture of radius \(a\) at \(k_{0}a=2.5\) claims dominant azimuthal content at \(\ell=10\). Determine the
maximum supported order \(\ell_{\mathrm{supp}}\) at \(\varepsilon=10^{-2}\) and test the claim.

\paragraph{Setup.}
Use Eq.~\eqref{eq:ch4.ell-max-scaling} and tail bound Eq.~\eqref{eq:ch4.high-ell-decay} on \(B_{a}\) \cite{hansen1988,jackson1999,stratton1941}.
VSH listing from Section~\ref{sec:ch34} \cite{hansen1988}.

\begin{figure}[t]
\centering
\ChOneFigBlock{%
\begin{tikzpicture}[ch1fig, node distance=7mm and 9mm]
  \node[ch1box] (ap) {Aperture $B_{a}$};
  \node[ch1box, right=11mm of ap] (ell) {$\ell_{\mathrm{supp}}$};
  \node[ch1box, right=11mm of ell] (tail) {Eq.~\eqref{eq:ch4.high-ell-decay}};
  \node[ch1box, below=11mm of ell] (claim) {claim $\ell=10$};
  \draw[ch1flow] (ap) -- (ell) -- (tail);
  \draw[ch1flow] (claim) -- (ell);
\end{tikzpicture}%
}
\caption{Example~\ref{sec:ch4101}: electrical size sets \(\ell_{\mathrm{supp}}\) before azimuthal claims.}
\label{fig:ch4.ex-ell-max-aperture}
\end{figure}

\textbf{Step 1.} Eq.~\eqref{eq:ch4.ell-max-scaling} gives \(\ell_{\mathrm{supp}}(10^{-2},2.5)\sim\mathcal{O}(6)\)--\(\mathcal{O}(8)\)
\cite{jackson1999,hansen1988,stratton1941}.

\textbf{Step 2.} Eq.~\eqref{eq:ch4.high-ell-decay} at \(\ell=10\), \(k_{0}a=2.5\) suppresses \(|\widehat{c}_{10,m}^{(\sigma)}|\) by
\(\exp(-10/2.5)=\exp(-4)\) relative to low-\(\ell\) sectors \cite{hansen1988,jackson1999}.

\textbf{Step 3.} Dominant \(\ell=10\) content fails \(\ell_{\mathrm{supp}}\) and tail audits: the claim is not in
\(\mathcal{S}_{\mathrm{real}}(\mathcal{C})\) on the declared aperture \cite{devaney2004,harrington2001}.

\textbf{Physical meaning.} Maximum OAM order on a compact aperture is a tail inequality, not a design preference \cite{hansen1988}.


\subsection{Construction of non-accessible eigenfields}
\label{sec:ch4102}
% TOC tag: [FIG+PROB][USE]

\textbf{Problem.} Construct a coefficient vector \(\mathbf{c}_{\mathrm{na}}\) that lies in \(\mathcal{U}_{\mathrm{unreach}}\) of
Eq.~\eqref{eq:ch4.unreachable-modes} for a centered dipole feed on \(B_{a}\) at \(k_{0}a=1\).

\begin{figure}[t]
\centering
\ChOneFigBlock{%
\begin{tikzpicture}[ch1fig, node distance=7mm and 9mm]
  \node[ch1box] (feed) {Dipole feed};
  \node[ch1box, right=11mm of feed] (gamma) {$\boldsymbol{\Gamma}_\omega$};
  \node[ch1box, right=11mm of gamma] (c) {$\mathbf{c}$};
  \node[ch1box, below=11mm of gamma] (null) {$\mathcal{U}_{\mathrm{unreach}}$};
  \draw[ch1flow] (feed) -- (gamma) -- (c);
  \draw[ch1flow] (null) -- (gamma);
\end{tikzpicture}%
}
\caption{Example~\ref{sec:ch4102}: unreachable directions are synthesis-null, not low coupling alone.}
\label{fig:ch4.ex-unreachable-modes}
\end{figure}

\textbf{Step 1.} For centered \(\Omega_{s}\), symmetry nulls odd-\(m\) and high-\(\ell\) sectors in Theorem~\ref{thm:ch4.symmetry-coupling-nulls}
\cite{hansen1988,jackson1999,stratton1941}.

\textbf{Step 2.} Choose \(\widehat{c}_{5,3}^{(\mathbf{M})}=1\) with all other entries zero \cite{harrington2001}.

\textbf{Step 3.} \(\mathbf{c}_{\mathrm{na}}\in\mathcal{U}_{\mathrm{unreach}}\): \(\delta_{\mathrm{syn}}>0\) despite unit coefficient norm
\cite{devaney2004,stratton1941}.

\textbf{Physical meaning.} Non-accessible eigenfields are symmetry- or rank-null directions, not ``weakly coupled'' modes \cite{hansen1988}.


\subsection{Near-field versus far-field truncation}
\label{sec:ch4103}
% TOC tag: [FIG+PROB][USE]

\textbf{Problem.} A near-field scan fits large \(k_{t}>k_{0}\) content. Estimate far-zone \(\mathbf{c}\) after propagating-sector projection.

\begin{figure}[t]
\centering
\ChOneFigBlock{%
\begin{tikzpicture}[ch1fig, node distance=7mm and 8mm]
  \node[ch1box] (nf) {Near scan};
  \node[ch1box, right=10mm of nf] (prop) {$\Pi_{\mathrm{prop}}$};
  \node[ch1box, right=10mm of prop] (ff) {Far $\mathbf{c}$};
  \node[ch1box, below=10mm of prop] (ev) {$\mathcal{K}_{\mathrm{ev}}$};
  \draw[ch1flow] (nf) -- (prop) -- (ff);
  \draw[ch1flow] (ev) -- (prop);
\end{tikzpicture}%
}
\caption{Example~\ref{sec:ch4103}: evanescent scan energy is removed before far-zone listing.}
\label{fig:ch4.ex-nf-ff-truncation}
\end{figure}

\textbf{Step 1.} Split scan spectrum into \(\mathcal{E}_{\mathrm{prop}}\) and \(\mathcal{E}_{\mathrm{ev}}\) via
Eq.~\eqref{eq:ch4.prop-ev-energy-fractions} \cite{stratton1941,jackson1999}.

\textbf{Step 2.} Apply Theorem~\ref{thm:ch4.prop-sector-equivalence}: only \(\widetilde{\mathbf{A}}_{\mathrm{syn}}\) enters \(\mathbf{c}\)
\cite{harrington2001,devaney2004}.

\textbf{Step 3.} If \(\mathcal{E}_{\mathrm{ev}}\gg\mathcal{E}_{\mathrm{prop}}\), far-zone \(\mathbf{c}\) is low-rank despite rich near-field
\cite{stratton1941,hansen1988}.

\textbf{Physical meaning.} Near-field truncation excess is reactive certification, not far-zone synthesis rank \cite{harrington2001}.


\subsection{Mode-count scaling with geometry and frequency}
\label{sec:ch4104}
% TOC tag: [FIG+PROB][USE]

\textbf{Problem.} A square aperture of side \(W\) operates at \(k_{0}W=4\). Estimate \(N_{\mathrm{mode}}^{\mathrm{(ap)}}(10^{-2})\) and compare
to a \(32\)-element array.

\begin{figure}[t]
\centering
\ChOneFigBlock{%
\begin{tikzpicture}[ch1fig, node distance=7mm and 9mm]
  \node[ch1box] (w) {$W$, $k_{0}W=4$};
  \node[ch1box, right=11mm of w] (nm) {Eq.~\eqref{eq:ch4.mode-count-area}};
  \node[ch1box, right=11mm of nm] (deff) {$d_{\mathrm{eff}}$};
  \node[ch1box, below=11mm of nm] (el) {32 elements};
  \draw[ch1flow] (w) -- (nm) -- (deff);
  \draw[ch1flow] (el) -- (deff);
\end{tikzpicture}%
}
\caption{Example~\ref{sec:ch4104}: mode-count scaling bounds diversity below element count.}
\label{fig:ch4.ex-mode-count-scaling}
\end{figure}

\textbf{Step 1.} \(A_{a}=W^{2}\), \(k_{0}^{2}A_{a}=(k_{0}W)^{2}=16\) \cite{jackson1999,stratton1941}.

\textbf{Step 2.} Eq.~\eqref{eq:ch4.mode-count-area} gives \(N_{\mathrm{mode}}^{\mathrm{(ap)}}\sim\mathcal{O}(10^{1})\)--\(\mathcal{O}(10^{2})\) up to
\(C_{A}(\varepsilon)\) \cite{hansen1988,collin2001,harrington2001}.

\textbf{Step 3.} Thirty-two elements do not guarantee thirty-two independent beams: \(d_{\mathrm{eff}}(10^{-2})\le N_{\mathrm{mode}}^{\mathrm{(ap)}}\)
\cite{devaney2004,stratton1941}.

\textbf{Physical meaning.} Mode-count scaling laws bound array diversity before beamforming claims \cite{hansen1988}.
